\section{Explicit existence and uniqueness obstructions}
The following examples separate failure of existence from failure of
uniqueness, and show why the coefficient-ring condition alone is insufficient.

\begin{lemma}\label{lem:finite-conductor-coefficient-obstruction}
Let $R\subset S$ be a finite injective extension, with $R$ nonzero
and local. If the conductor $J=(R:S)$ has finite quotient $R/J$,
then the binary polynomial tensor map is injective exactly when
$S=R$.
\end{lemma}
\begin{proof}
The monic product $P$ over representatives of $R/J$ has
$P(R)\subset J$, and monic division gives
$\Omega(R,S)=E\oplus PS[X]$, where $E$ has degree below $\deg P$.
In particular $SP$ is a direct summand. Restriction to its tensor
square forces $S\otimes_RS\to S$ to be injective, hence an
isomorphism. A finite ring epimorphism is surjective
\cite[Lemma~10.107.6]{StacksRingEpimorphisms}, so injectivity of
$R\to S$ gives $R=S$. The converse is immediate.
\end{proof}

\begin{lemma}\label{lem:binary-mixed-difference-obstruction}
Let $(R,\mathfrak m)$ be a local domain. Suppose that $u_n\in R$
has the following property: for every finite family
$\mathcal H\subset\operatorname{Int}(R)$, all sufficiently large
$n$ satisfy $h(u_n)-h(0)\in\mathfrak m$ for every $h\in\mathcal H$.
If $g\in\operatorname{Int}(R)$ has $g(0)=0$ and
$g(u_n^2)\notin\mathfrak m^2$ for all $n$, then $g(XY)$ is not a
finite sum of products $h_i(X)k_i(Y)$ with integer-valued factors.
\end{lemma}
\begin{proof}
The mixed difference at $(0,0),(u_n,0),(0,u_n),(u_n,u_n)$ of such
an expansion would be
$\sum_i(h_i(u_n)-h_i(0))(k_i(u_n)-k_i(0))\in\mathfrak m^2$.
For $g(XY)$ the same difference is $g(u_n^2)$, a contradiction.
\end{proof}

\begin{example}\label{ex:all-finite-fields-two-branches}
For every finite field $k=\mathbb F_q$, let
\[
 T=k[t]_{k[t]\setminus((t)\cup(t-1))},\quad
 m=t(t-1),\quad R=k+mT,\quad K=k(t).
\]
Then $R$ is a one-dimensional Noetherian local domain with finite
normalization $T$. Put
\[
 a(X)=\frac{X^q-X}{m},\qquad g(X)=a(X)^q-a(X),\qquad
 u_n=t(t-1)^{n+1}.
\]
The polynomial $g$ is integer-valued on $R$, has degree $q^2$,
and $g(XY)$ has no multiplicative tensor expansion.
\end{example}
\begin{proof}
Let $A=k[m]_{(m)}$. The equation $t^2=t+m$ gives
$T=A\oplus At$: every inverted polynomial $f(t)$ has
$f(t)f(1-t)\in k[m]$ with nonzero constant term, and its
inverse is therefore already in $A+At$.
Also $R=A\oplus A(mt)$ with
$(mt)^2=m(mt)+m^3$. Hence $R$ is finite over the DVR $A$ and
is Noetherian of dimension one. Its maximal ideal is $mT$ and
its residue field is $k$. Its fraction field is $K$, and the
normal finite overring $T$ is its normalization.

For $r\in R$, $r^q-r\in mT$, so $a(r)\in T$; reduction at
both maximal ideals of $T$ gives $g(r)\in mT\subset R$.
Write $v_0,v_1$ for the two normalized valuations. One has
\[
 v_0(u_n)=1,\quad v_1(u_n)=n+1,\quad
 v_0(g(u_n^2))=1,\quad v_1(g(u_n^2))=2n+1.
\]
For the last two equalities, $u_n^2$ has positive value at both
branches, so the leading term of $x^q-x$ is $-x$; division by
$m$ subtracts one from the value, and the same observation applies
to $a^q-a$. For finitely many polynomials, every nonconstant term
of $h(u_n)-h(0)$ tends to zero at $v_1$, while the difference lies
in $R$. It therefore belongs to $(t-1)T\cap R=mT$ eventually.
Since $(mT)^2$ has $v_0$ at least two, the preceding lemma applies.
\end{proof}

\begin{example}\label{ex:two-dimensional-split-tensor-kernel}
Let $T=\mathbb F_2[x,y]_{(x,y)}$, $\mathfrak n=(x,y)T$, and
$R=\mathbb F_2+\mathfrak n^2$. This is a two-dimensional
Noetherian local domain with finite normalization $T$. For
$q(X)=X^2-X$,
\[
 \operatorname{Int}(R)=R\oplus RX\oplus q(X)T[X],\qquad
              (xq)\otimes q-q\otimes(xq)\ne0,
\]
although the displayed tensor has zero polynomial image.
\end{example}
\begin{proof}
Over $A=\mathbb F_2[x^2,y^2]_{(x^2,y^2)}$, the ring $T$ is finite
with generators $1,x,y,xy$, and $R$ is generated by
$1,xy,x^3,xy^2,x^2y,y^3$. Thus both are Noetherian, and integrality
gives dimension two. The maximal ideal of $R$ is $\mathfrak n^2$;
its fraction field equals that of $T$, since $x=x^3/x^2$ and
$y=y^3/y^2$. The normal ring $T$ is therefore the finite
normalization.

At each height-one prime of $T$, the image of $R$ is infinite,
so Vandermonde interpolation puts every integer-valued coefficient
in that localization. Normality gives coefficients in $T$.
Since $q(R)\subset\mathfrak n^2$ and this is a $T$-ideal, monic
division by $q$ reduces the value condition to a linear remainder;
evaluation at $0,1$ puts its coefficients in $R$. This proves the
decomposition. In its split summand $Tq\otimes_RTq$, reduction
modulo $\mathfrak n^2$ gives $D\otimes_{\mathbb F_2}D$, where
$D=\mathbb F_2[x,y]/(x,y)^2$. The tensor $x\otimes1-1\otimes x$
is nonzero in its nine-dimensional monomial basis and maps to
zero under multiplication. This detects the displayed tensor.
\end{proof}

\begin{lemma}\label{lem:derivation-detects-polynomial-torsion}
Let $R$ be a domain with fraction field $K$, let $V\subset K$ be
a DVR containing $R$, and let $D:K\to K$ be a derivation with
$D(R)\subset V$. Suppose $u_n\in R$ satisfies $v(u_n)\to\infty$
and $D(u_n)=1$. Then $F'(0)\in V$ for every
$F\in\operatorname{Int}(R)$, and
\[
       \beta(F,G)=D(F'(0))G'(0)+V\quad\text{in }K/V
\]
is balanced over $R$. If $P\in R[X]$, $cP\in\operatorname{Int}(R)$,
and $D(c)P'(0)^2\notin V$, then
$(cP)\otimes P-P\otimes(cP)$ is nonzero and has zero polynomial
image.
\end{lemma}
\begin{proof}
The chain rule gives
$D(F(u_n))=(DF)(u_n)+F'(u_n)$. Its left side lies in $V$;
the constant coefficient of $DF$ is $D(F(0))\in V$, and its other
terms eventually have nonnegative valuation. Thus $F'(u_n)\in V$,
and the same estimate for $F'(u_n)-F'(0)$ gives $F'(0)\in V$.
For $r\in R$, the balancing error is
$D(r)F'(0)G'(0)\in V$. Finally evaluation on the displayed tensor
cancels the two $D(P'(0))$ terms and leaves
$D(c)P'(0)^2+V\ne0$. Its polynomial image cancels identically.
\end{proof}

\begin{example}\label{ex:akizuki-reid-polynomial-kernel}
Let $A_0$ be a DVR with parameter $t$ and finite residue field
$\mathbb F_q$. Choose units $a_i\in A_0$ and integers
$n_0=0$, $n_{i+1}\ge2n_i+2$, such that
$z=\sum_i a_it^{n_i}\in\widehat A_0$ is transcendental over
$\operatorname{Frac}(A_0)$. Put
\[
 z_i=\frac{z-\sum_{j<i}a_jt^{n_j}}{t^{n_i}},\quad
 C=A_0[t(z-a_0),(z_i-a_i)^2:i\ge0],\quad
 M=(t,t(z-a_0))C,\quad R=C_M.
\]
Then $R$ is a one-dimensional Noetherian local domain whose
normalization is a DVR with parameter $t$ and the same residue
field, but is not finite over $R$. For $c=z-a_0$ and $P=X^q-X$,
\[
 cP\in\operatorname{Int}(R),\qquad
              (cP)\otimes P-P\otimes(cP)\ne0
\]
with zero polynomial image. This holds in mixed characteristic
as well as positive characteristic.
\end{example}
\begin{proof}
The ring assertions are Reid's two theorems on page~2, under
precisely these hypotheses; their proofs occupy pages~2--4
\cite{Reid1995Akizuki}. We give the obstruction and a separate
nonfiniteness argument. On $K=\operatorname{Frac}(A_0)(z)$ set
$D=t^{-1}\partial/\partial z$. It kills $A_0$, sends
$t(z-a_0)$ to $1$, and
\[
              D((z_i-a_i)^2)=2t^{-n_i-1}(z_i-a_i)
\]
has value at least $n_{i+1}-2n_i-1\ge1$ (or is zero).
The quotient rule gives $D(R)\subset V$. The elements
$u_i=t(z-\sum_{j\le i}a_jt^{n_j})\in R$ have derivative $1$
and value $n_{i+1}+1\to\infty$.

Writing $u=t(z-a_0)$, one has $ct=u$ and $cu=t(z-a_0)^2\in C$,
so $cM R\subset R$. Since $P(R)\subset MR$, $cP$ is
integer-valued. The preceding lemma detects its tensor by
$D(c)P'(0)^2=1/t\notin V$.
If $V$ were finite over $R$, the derivatives of a finite module
generating set would bound the valuations of all $D(V)$ below,
because $D(R)\subset V$. But $z_i\in V$ and
$D(z_i)=t^{-n_i-1}$, a contradiction.

An entirely explicit characteristic-two instance is
$A_0=\mathbb F_2[t]_{(t)}$, $a_i=1$, $n_i=(i+1)!$ for $i\ge1$.
The gap condition holds. To verify transcendence, a hypothetical
relation $P(t,z)=0$ of $z$-degree $d$ and coefficient degrees at
most $h$ would make its value on the $N$th truncation have
valuation at least $n_{N+1}$ and degree at most $h+dn_N$.
For large $N$ it must vanish identically. Infinitely many distinct
truncations cannot be roots of the same nonzero polynomial over
$\mathbb F_2(t)$, giving the required contradiction.
\end{proof}

\begin{proposition}\label{prop:coefficient-basechange-not-sufficient}
For any $R\subset S$, let $C$ be its effective coefficient ring,
$Q=C/R$, and
\[
 N(R,Q)=\{(q_i)\in\bigoplus_{i\ge0}Q:
                         \sum_i r^iq_i=0\text{ for every }r\in R\}.
\]
Binary tensor injectivity implies
\[
 C\otimes_R\Omega(R,S)\simeq C[X],\qquad
 C\otimes_RN(R,Q)=0.
\]
These two conditions are equivalent even without tensor injectivity.
They are not sufficient, even when $C$ is a flat epic localization
of a domain. In that scope a nonzero tensor can be killed by a
monic polynomial in one variable.
\end{proposition}
\begin{proof}
The coefficient reduction sequence is
$0\to R[X]\to\Omega\to N(R,Q)\to0$: lift a finite vector of
coefficients and use its defining equations for surjectivity.
After tensoring by $C$, the natural map $C[X]\to C\otimes_R\Omega$
is split by multiplication. Its cokernel is $C\otimes_RN(R,Q)$,
proving the equivalence.
If $z=\sum c_i\otimes F_i$ maps to zero in $C[X]$, choose a common
monic $H\in R[X]$ with every $c_iH\in\Omega$, using the monic
coefficient witnesses and their product. The tensor
$\sum c_iH\otimes F_i$ has zero polynomial image. Under tensor
injectivity it is zero, and extracting its first leading
coefficient gives $z=0$.

For the counterexample let $R=k+tV$, $V=l[[t]]$, $K=l((t))$,
where $k=\mathbb F_q\subset l$ and $l$ is infinite. Put $P=X^q-X$.
The valuation-conductor decomposition gives
\[
 \operatorname{Int}(R)=\bigoplus_{i<q}RX^i\ \oplus\!
 \bigoplus_{i<q,j\ge1}t^{1-j}VX^iP^j.
\]
Every element of $K$ occurs as a leading coefficient, so $C=K$
and $K\otimes_R\operatorname{Int}(R)=K[X]$.
Choose $\beta\in l\setminus k$. The tensor
$z=\beta P\otimes P-P\otimes\beta P$ has zero polynomial
image but projects to the nonzero element
$\beta\otimes1-1\otimes\beta\in l\otimes_kl$.
Nevertheless $P^2/t$ is integer-valued and $t,t\beta\in R$, so
\[
 \beta P^2\otimes P=(P^2/t)\otimes(t\beta P)
                         =P^2\otimes\beta P.
\]
Thus $P(X)z=0$.
More generally, for $R=k+t^aV$, every proper rational coefficient
overring fails uniqueness by the valuation-conductor theorem:
an overring inside $V$ lies in $(R:_S t^aV)$, and an overring
outside $V$ is $K$. Indeed, a sufficiently large power of an
element of negative value has inverse in $t^{a+1}V$; multiplying
the conductor by that power puts both $V$ and $t^{-1}$ in the
overring. For $F=P^2/t^a$, the split coefficient module
$t^{-a}V$ detects
$(\beta F)\otimes F-F\otimes(\beta F)\ne0$ after reduction to
$l\otimes_kl$. These rings are non-Noetherian, since
$t^aV/t^{2a}V$ has infinite dimension over $k$.
\end{proof}

\begin{corollary}\label{cor:flat-cusp-nonzero-tor-obstruction}
For the cusp with $a=1$ in the preceding proof, the value module is
\[
 E=\bigoplus_{i<q}(K/R)X^i\ \oplus\!
   \bigoplus_{i<q,j\ge1}(K/t^{1-j}V)X^iP^j.
\]
The universal obstruction $K_2=\operatorname{Tor}_2^R(E,E)$ is
nonzero. Its nonzero class is the double connecting lift of
$\beta\otimes1-1\otimes\beta\in V\otimes_RV$ from the exact
sequence $0\to V\to K\to K/V\to0$.
\end{corollary}
\begin{proof}
The displayed description follows by quotienting the split monic
coefficient modules. The $P$-summand is $K/V$. The indicated tensor
is nonzero by residue reduction and dies in $K\otimes_RV=K$.
Since $K$ is flat, two long exact Tor sequences identify that
kernel with $\operatorname{Tor}_2^R(K/V,K/V)$. The split inclusion
of $K/V$ into $E$ preserves this nonzero class. The sign is the
connecting convention already fixed in the universal theorem;
nonvanishing is independent of that sign.
\end{proof}

\begin{corollary}\label{cor:full-bivariate-flat-obstruction}
Under the flat epic coefficient hypothesis, evaluation modulo
$C/R$ induces a map
$\lambda:\operatorname{Tor}_1^R(E,E)\to\operatorname{App}(R^2,C/R)$
through the chosen connecting map. There is an exact sequence
\[
 0\to\operatorname{Tor}_2^R(E,E)\to\Omega^{\otimes_R2}
                \to\Omega_2\to\ker\lambda\to0.
\]
Thus the full bivariate tensor map is an isomorphism exactly when
$\operatorname{Tor}_2^R(E,E)=0$ and $\lambda$ is injective.
\end{corollary}
\begin{proof}
Evaluation modulo $C/R$ vanishes on the tensor image, so it factors
uniquely through the cokernel and defines $\lambda$. A class of
$\ker\lambda$ lifts through the surjective connecting map to a
polynomial whose values are in $R$, hence to $\Omega_2$.
The stated sequence and criterion follow. This retains the
distinction between the two specified substitution tests and
surjectivity for every bivariate integer-valued polynomial.
\end{proof}
